\section{Agent 是新物种}

上一章讨论智能衡量，本章进入 Agent。广密把 Agent 称为“新物种”，是因为它不再只是语言输出系统，而是能感知环境、保持任务状态、调用工具、执行行动、读取反馈并继续推进目标的系统。越接近 AGI，Agent 的变化可能越像阶段跃迁，而不是线性增强。

\begin{figure}[H]
\centering
\includegraphics[width=0.92\textwidth]{agent-three-capabilities.png}
\caption{Agent 三个关键能力：Long context reasoning、Tool use、Instruction following 共同支撑新物种。自制概念图，依据 00:48:03--00:55:49 对谈内容整理。}
\end{figure}

\begin{knowledgebox}{读图：三个能力缺一不可}
Long Context 让 Agent 保持长任务状态，Tool Use 让它改变外部世界，Instruction Following 让它不偏离用户目标。Memory、Feedback 和 Autonomy 是进一步升级：Agent 不只执行当前命令，还能积累经验和主动探索。
\end{knowledgebox}

\subsection{Long context reasoning：长任务状态}

本节先看长上下文推理。Agent 任务往往不是一句问答，而是多步骤、多文件、多工具、多轮错误恢复。模型需要记住目标、约束、已尝试路径、失败原因、当前产物和下一步。长上下文只是最低门槛，真正难点是对上下文做选择、压缩和更新。

\begin{warningbox}{长上下文不是万能记忆}
把所有内容塞进窗口，会带来成本、噪声和注意力稀释。好的 Agent 需要结构化任务状态、外部文件系统、检索、摘要和长期记忆，而不是无限堆 token。
\end{warningbox}

\subsection{Tool use：从语言到行动}

本节看工具使用。工具调用让模型能搜索、浏览网页、读写文件、运行代码、调用 API、操作 UI。它把模型从“产生建议”推向“执行任务”。但工具使用也带来安全、权限、错误恢复和可审计性问题。越强的 Agent，越需要明确的边界和可回滚机制。

\begin{importantbox}{Agent 产品的基本契约}
用户给目标，Agent 拆任务并调用工具；系统必须提供权限控制、进度可见性、错误解释、人工接管和结果验证。没有这些，Agent 越强，风险越大。
\end{importantbox}

\subsection{Instruction following：复杂目标的服从}

本节解释第三个能力。Agent 不只是会用工具，还要在长程任务中持续服从用户目标：哪些事情不能做，哪些资源不能用，什么结果才算完成，什么时候需要确认，如何处理冲突指令。复杂 instruction following 是 Agent 从玩具 demo 走向生产系统的关键。

\noindent\begin{tabularx}{\textwidth}{p{0.20\textwidth}p{0.34\textwidth}X}
\toprule
能力 & 失败模式 & 产品要求 \\
\midrule
长上下文推理 & 忘记目标、重复尝试、误用旧信息 & 任务状态和外部记忆。 \\
工具使用 & 调错工具、破坏文件、权限越界 & 沙盒、确认、审计和回滚。 \\
指令遵循 & 被中途信息带偏、忽略约束 & 明确优先级和冲突处理。 \\
反馈学习 & 同错再犯、无法迁移经验 & 保存失败、复盘和评测闭环。 \\
\bottomrule
\end{tabularx}

\subsection{本章小结}

Agent 是新物种，不是因为它有一个新名字，而是因为它把语言模型接入环境、工具、记忆和行动闭环。真正的竞争会从“回答质量”转向“任务完成质量”。

